\begin{lemma}
\label{lemma-alternating-cech-to-cohomology}
Let $S$ be a scheme. Let $f : U \to X$ be a surjective \'etale morphism
of algebraic spaces over $S$. Let $\mathcal{F}$ be an object of
$\textit{Ab}(X_\etale)$. There exists a canonical map
$$
\check{\mathcal{C}}^\bullet_{alt}(f, \mathcal{F})
\longrightarrow
R\Gamma(X, \mathcal{F})
$$
in $D(\textit{Ab})$. Moreover, there is a spectral sequence with $E_1$-page
$$
E_1^{p, q} =
\Ext_{\textit{Ab}(X_\etale)}^q(K^p, \mathcal{F})
$$
converging to $H^{p + q}(X, \mathcal{F})$ where
$K^p = \wedge^{p + 1}f_!\underline{\mathbf{Z}}$.
\end{lemma}

\begin{proof}
Recall that we have the quasi-isomorphism
$K^\bullet \to \underline{\mathbf{Z}}[0]$, see
(\ref{equation-quasi-isomorphism}).
Choose an injective resolution $\mathcal{F} \to \mathcal{I}^\bullet$
in $\textit{Ab}(X_\etale)$. Consider the double complex
$\Hom(K^\bullet, \mathcal{I}^\bullet)$ with terms
$\Hom(K^p, \mathcal{I}^q)$. The differential
$d_1^{p, q} : A^{p, q} \to A^{p + 1, q}$
is the one coming from the differential $K^{p + 1} \to K^p$
and the differential $d_2^{p, q} : A^{p, q} \to A^{p, q + 1}$ is the
one coming from the differential
$\mathcal{I}^q \to \mathcal{I}^{q + 1}$.
Denote $\text{Tot}(\Hom(K^\bullet, \mathcal{I}^\bullet))$
the associated total complex, see
Homology, Section \ref{homology-section-double-complexes}.
We will use the two spectral
sequences $({}'E_r, {}'d_r)$ and $({}''E_r, {}''d_r)$
associated to this double complex, see
Homology, Section \ref{homology-section-double-complex}.

\medskip\noindent
Because $K^\bullet$ is a resolution of $\underline{\mathbf{Z}}$
we see that the complexes
$$
\Hom(K^\bullet, \mathcal{I}^q) :
\Hom(K^0, \mathcal{I}^q) \to
\Hom(K^1, \mathcal{I}^q) \to
\Hom(K^2, \mathcal{I}^q) \to \ldots
$$
are acyclic in positive degrees and have $H^0$ equal to
$\Gamma(X, \mathcal{I}^q)$. Hence by
Homology, Lemma \ref{homology-lemma-double-complex-gives-resolution}
the natural map
$$
\mathcal{I}^\bullet(X) \longrightarrow
\text{Tot}(\Hom(K^\bullet, \mathcal{I}^\bullet))
$$
is a quasi-isomorphism of complexes of abelian groups. In particular
we conclude that
$H^n(\text{Tot}(\Hom(K^\bullet, \mathcal{I}^\bullet))) = H^n(X, \mathcal{F})$.

\medskip\noindent
The map $\check{\mathcal{C}}^\bullet_{alt}(f, \mathcal{F}) \to
R\Gamma(X, \mathcal{F})$ of the lemma is the composition of
$\check{\mathcal{C}}^\bullet_{alt}(f, \mathcal{F}) \to
\text{Tot}(\Hom(K^\bullet, \mathcal{I}^\bullet))$
with the inverse of the displayed quasi-isomorphism.

\medskip\noindent
Finally, consider the spectral sequence $({}'E_r, {}'d_r)$.
We have
$$
E_1^{p, q} = q\text{th cohomology of }
\Hom(K^p, \mathcal{I}^0) \to
\Hom(K^p, \mathcal{I}^1) \to
\Hom(K^p, \mathcal{I}^2) \to \ldots
$$
This proves the lemma.
\end{proof}

\noindent
It follows from the lemma that it is important to understand the
ext groups $\Ext_{\textit{Ab}(X_\etale)}(K^p, \mathcal{F})$,
i.e., the right derived functors of
$\mathcal{F} \mapsto \Hom(K^p, \mathcal{F})$.

\begin{lemma}
\label{lemma-compute}
Let $S$ be a scheme. Let $f : U \to X$ be a surjective, \'etale, and separated
morphism of algebraic spaces over $S$. For $p \geq 0$ set
$$
W_p = U \times_X \ldots \times_X U \setminus \text{all diagonals}
$$
where the fibre product has $p + 1$ factors.
There is a free action of $S_{p + 1}$ on $W_p$ over $X$ and
$$
\Hom(K^p, \mathcal{F}) = S_{p + 1}\text{-anti-invariant elements of }
\mathcal{F}(W_p)
$$
functorially in $\mathcal{F}$ where
$K^p = \wedge^{p + 1}f_!\underline{\mathbf{Z}}$.
\end{lemma}

\begin{proof}
Because $U \to X$ is separated the diagonal $U \to U \times_X U$ is a
closed immersion. Since $U \to X$ is \'etale the diagonal
$U \to U \times_X U$ is an open immersion, see
Morphisms of Spaces, Lemmas
\ref{spaces-morphisms-lemma-etale-unramified} and
\ref{spaces-morphisms-lemma-diagonal-unramified-morphism}.
Hence $W_p$ is an open and closed subspace of
$U^{p + 1} = U \times_X \ldots \times_X U$. The action of $S_{p + 1}$
on $W_p$ is free as we've thrown out the fixed points of the action.
By
Lemma \ref{lemma-product-is-tensor-product}
we see that
$$
(f_!\underline{\mathbf{Z}})^{\otimes p + 1} =
f^{p + 1}_!\underline{\mathbf{Z}} = (W_p \to X)_!\underline{\mathbf{Z}}
\oplus Rest
$$
where $f^{p + 1} : U^{p + 1} \to X$ is the structure morphism.
Looking at stalks over a geometric point $\overline{x}$ of $X$
we see that
$$
\left(
\bigoplus\nolimits_{\overline{u} \mapsto \overline{x}} \mathbf{Z}
\right)^{\otimes p + 1}
\longrightarrow
(W_p \to X)_!\underline{\mathbf{Z}}_{\overline{x}}
$$
is the quotient whose kernel is generated by all tensors
$1_{\overline{u}_0} \otimes \ldots \otimes 1_{\overline{u}_p}$
where $\overline{u}_i = \overline{u}_j$ for some $i \not = j$.
Thus the quotient map
$$
(f_!\underline{\mathbf{Z}})^{\otimes p + 1}
\longrightarrow
\wedge^{p + 1}f_!\underline{\mathbf{Z}}
$$
factors through $(W_p \to X)_!\underline{\mathbf{Z}}$, i.e., we get
$$
(f_!\underline{\mathbf{Z}})^{\otimes p + 1}
\longrightarrow
(W_p \to X)_!\underline{\mathbf{Z}}
\longrightarrow
\wedge^{p + 1}f_!\underline{\mathbf{Z}}
$$
This already proves that $\Hom(K^p, \mathcal{F})$ is (functorially) a
subgroup of
$$
\Hom((W_p \to X)_!\underline{\mathbf{Z}}, \mathcal{F}) = \mathcal{F}(W_p)
$$
To identify it with the $S_{p + 1}$-anti-invariants we have to prove that
the surjection $(W_p \to X)_!\underline{\mathbf{Z}}
\to \wedge^{p + 1}f_!\underline{\mathbf{Z}}$ is the maximal
$S_{p + 1}$-anti-invariant quotient. In other words, we have to show that
$\wedge^{p + 1}f_!\underline{\mathbf{Z}}$ is the quotient of
$(W_p \to X)_!\underline{\mathbf{Z}}$ by the subsheaf generated by
the local sections $s - \text{sign}(\sigma)\sigma(s)$ where $s$ is
a local section of $(W_p \to X)_!\underline{\mathbf{Z}}$.
This can be checked on the stalks, where it is clear.
\end{proof}

\begin{lemma}
\label{lemma-twist}
Let $S$ be a scheme. Let $W$ be an algebraic space over $S$.
Let $G$ be a finite group acting freely on $W$.
Let $U = W/G$, see
Properties of Spaces, Lemma \ref{spaces-properties-lemma-quotient}.
Let $\chi : G \to \{+1, -1\}$ be a character.
Then there exists a rank 1 locally free sheaf of $\mathbf{Z}$-modules
$\underline{\mathbf{Z}}(\chi)$ on $U_\etale$ such that for every
abelian sheaf $\mathcal{F}$ on $U_\etale$ we have
$$
H^0(W, \mathcal{F}|_W)^\chi =
H^0(U, \mathcal{F} \otimes_{\mathbf{Z}} \underline{\mathbf{Z}}(\chi))
$$
\end{lemma}

\begin{proof}
The quotient morphism $q : W \to U$ is a $G$-torsor, i.e., there exists
a surjective \'etale morphism $U' \to U$ such that
$W \times_U U' = \coprod_{g \in G} U'$ as spaces with $G$-action over $U'$.
(Namely, $U' = W$ works.) Hence $q_*\underline{\mathbf{Z}}$ is a finite
locally free $\mathbf{Z}$-module with an action of $G$. For any
geometric point $\overline{u}$ of $U$, then we get $G$-equivariant
isomorphisms
$$
(q_*\underline{\mathbf{Z}})_{\overline{u}}
= \bigoplus\nolimits_{\overline{w} \mapsto \overline{u}} \mathbf{Z}
= \bigoplus\nolimits_{g \in G} \mathbf{Z} = \mathbf{Z}[G]
$$
where the second $=$ uses a geometric point
$\overline{w}_0$ lying over $\overline{u}$ and
maps the summand corresponding to $g \in G$ to the summand
corresponding to $g(\overline{w}_0)$. We have
$$
H^0(W, \mathcal{F}|_W) =
H^0(U, \mathcal{F} \otimes_\mathbf{Z} q_*\underline{\mathbf{Z}})
$$
because
$q_*\mathcal{F}|_W = \mathcal{F} \otimes_\mathbf{Z} q_*\underline{\mathbf{Z}}$
as one can check by restricting to $U'$. Let
$$
\underline{\mathbf{Z}}(\chi) =
(q_*\underline{\mathbf{Z}})^\chi \subset
q_*\underline{\mathbf{Z}}
$$
be the subsheaf of sections that transform according to $\chi$. For
any geometric point $\overline{u}$ of $U$ we have
$$
\underline{\mathbf{Z}}(\chi)_{\overline{u}} =
\mathbf{Z} \cdot \sum\nolimits_g \chi(g) g
\subset
\mathbf{Z}[G] = (q_*\underline{\mathbf{Z}})_{\overline{u}}
$$
It follows that $\underline{\mathbf{Z}}(\chi)$ is locally free of
rank 1 (more precisely, this should be checked after restricting to $U'$).
Note that for any $\mathbf{Z}$-module $M$ the $\chi$-semi-invariants
of $M[G]$ are the elements of the form $m \cdot \sum\nolimits_g \chi(g) g$.
Thus we see that for any abelian sheaf $\mathcal{F}$ on $U$ we have
$$
\left(\mathcal{F} \otimes_\mathbf{Z} q_*\underline{\mathbf{Z}}\right)^\chi
=
\mathcal{F} \otimes_\mathbf{Z} \underline{\mathbf{Z}}(\chi)
$$
because we have equality at all stalks. The result of the lemma follows by
taking global sections.
\end{proof}
